%%=============================================================================
%% genAI
%%=============================================================================

\chapter*{\IfLanguageName{dutch}{Gebruik Generatieve AI}{Use of Generative AI}}%
\label{ch:genAI}

Doorheen deze bachelorproef en mijn stage heb ik gebruikgemaakt van de gratis plannen van zowel Gemini, Claude en Copilot. Bij Gemini heb ik vooral het `thinking` model gebruikt wanneer beschikbaar en bij Claude zijn de `adaptive thinking` optie met het sonnet 4.6 model. Alle drie de toepassingen heb ik gebruikt voor code suggesties en de code na te lezen en verbeteren. Dit is gedaan in stukken en steeds nog verwerkt door mij vooraleer het werd gecommit op de repository. Enkel het `wireviz\_generator` bestand van de python package voor de \gls{MES} is \textit{gevibecode} door de korte tijd die nog over was en de complexiteit van Wireviz. Omdat ik deze bachelorproef ook geschreven heb in Latex in tegenstelling tot Word heb ik ook Claude gebruikt om het voorblad te converteren in een latex document class. Deze heb ik doorgestuurd naar mijn interne promotor voor toekomstig gebruik. Ook zal ik mijn volledige repository doorsturen wanneer deze is opgekuist.

Daarnaast heb ik zowel Gemini als Claude gebruikt om bronnen op te zoeken die mij konden helpen in meer informatie op te doen over mijn onderwerp en mij, wanneer ik extra uitleg nodig had, de informatie uitlegde. Als laatste heb ik beide gebruikt om mijn geschreven teksten na te lezen op slechte grammatica en fouten die erin beland zouden zijn, om de zinsbouw te verbeteren en om te controleren of de hoofdstukken in lijn liggen met het vademecum.